\chapter{对数积分与指数积分：\texorpdfstring{$\Li$}{Li}、\texorpdfstring{$\li$}{li} 与 \texorpdfstring{$\Ei$}{Ei}}\label{ch:logint_Ei}
\label{ch:Li_li_Ei}

% ============================================================
\section*{本章导读}
\addcontentsline{toc}{section}{本章导读}
% ============================================================

\(\li\) 为从 $0$ 取主值的对数积分，可延拓为复函数 \(\li(z)\)；
\(\Li\) 为自 $2$ 起算的实偏移版，与 \(\li\) 仅差常数 \(\li(2)\)。
\(\Ei\) 是 \(\li\) 经对数换元后的等价形式。
素数定理用实 \(\Li\)；复点上用 \(\li(z)\)；数值实现多用 \(\Ei\)。

% ============================================================
\section{实对数积分 \texorpdfstring{$\li(x)$}{li(x)}}
\label{sec:Li-li-summary}
% ============================================================

\begin{definition}{对数积分（实轴）}{}
\label{def:li-real}
对 $x>1$，
\begin{equation}
  \li(x)
  \;=\;
  \mathrm{P.V.}\int_0^x \frac{\mathrm{d}t}{\log t}
  \;=\;
  \lim_{\varepsilon\to 0^+}
  \Biggl(
    \int_0^{1-\varepsilon}\frac{\mathrm{d}t}{\log t}
    +\int_{1+\varepsilon}^{x}\frac{\mathrm{d}t}{\log t}
  \Biggr).
  \label{eq:li-def-ch2}
\end{equation}
\end{definition}

被积函数在 $t=1$ 有奇点，故取 Cauchy 主值。对 $x>1$，$\li(x)$ 为实数。记
\begin{equation}
  \li(2)
  \;=\;
  \mathrm{P.V.}\int_0^2\frac{\mathrm{d}t}{\log t}
  \;\approx\;
  1.04516.
  \label{eq:li-of-two}
\end{equation}

\subsection{主值在 \texorpdfstring{$t=1$}{t=1} 附近的收敛性}

在 $t=1$ 附近令 $t=e^{u}$（$u\to 0$），则
\[
  \frac{\mathrm{d}t}{\log t}
  =\frac{e^{u}}{u}\,\mathrm{d}u
  =\Bigl(\frac{1}{u}+1+O(u)\Bigr)\mathrm{d}u.
\]
奇部 $1/u$ 关于 $u=0$ 为奇函数，主值贡献为零；$1+O(u)$ 可积。
故~\eqref{eq:li-def-ch2} 跨过 $t=1$ 的极限存在（有限）。

\begin{figure}[htbp]
\centering
\begin{tikzpicture}[
  >=Stealth,
  font=\small,
  every node/.style={inner sep=1pt}
]
  \draw[->] (-0.4,0) -- (8.2,0) node[right] {$t$（实轴）};
  \draw[->] (0,-1.1) -- (0,1.4);
  \foreach \x/\lab in {0/0, 1.8/1, 3.2/2, 6.5/x}
    \draw (\x,0.08) -- (\x,-0.08) node[below] {$\lab$};
  \fill[red!70!black] (1.8,0) circle (1.6pt);
  \node[red!70!black, above] at (2.3,0.65) {奇点 $t=1$};% 相对原位右移0.5cm、上移0.5cm
  \draw[blue!75!black, very thick, ->]
    (0.08,0) -- (1.35,0);
  \node[blue!75!black, above] at (0.7,0.05) {$\int_0^{1-\varepsilon}$};
  \draw[blue!75!black, very thick]
    (1.35,0) arc (180:0:0.45);
  \node[blue!75!black, above] at (2.3,1.05) {主值挖去};% 相对原位右移0.5cm、上移0.5cm
  \draw[blue!75!black, very thick, ->]
    (2.25,0) -- (6.45,0);
  \node[blue!75!black, above] at (4.3,0.05) {$\int_{1+\varepsilon}^{x}$};
  \draw[green!50!black, thick, dashed, ->]
    (3.2, -0.55) -- (6.45, -0.55);
  \node[green!50!black, below] at (4.8,-0.55) {$\Li$：自 $2$ 至 $x$（不过 $t=1$）};
  \node[blue!75!black, below left] at (7.6,-0.15) {$\li(x)$};
\end{tikzpicture}
\caption{实轴上的 $\li(x)$（$x>1$）：在 $t=1$ 处取 Cauchy 主值。
虚线为 $\Li(x)=\int_2^x$，不过奇点。}
\label{fig:li-real-path}
\end{figure}

$x\to+\infty$ 时的渐近与第~\ref{ch:prime_distribution}~章对 $\Li$ 的分部积分同型：
\begin{equation}
  \li(x)
  \;\sim\;
  \frac{x}{\log x}
  +\frac{x}{(\log x)^{2}}
  +\cdots
  \qquad(x\to+\infty).
  \label{eq:li-asymp-real}
\end{equation}

% ============================================================
\section{偏移对数积分 \texorpdfstring{$\Li$}{Li}：实函数及其与 \texorpdfstring{$\li$}{li} 的关系}
\label{sec:Li}
% ============================================================

\begin{definition}{偏移对数积分（实函数）}{}
\label{def:Li-real}
$\Li$ 是定义在实半轴上的实值函数
\begin{equation}
  \Li \colon (1,+\infty)\to\mathbb{R},
  \qquad
  \Li(x)
  \;=\;
  \int_2^x \frac{\mathrm{d}t}{\log t}.
  \label{eq:Li-def-ch2}
\end{equation}
下限取 $2$，不过 $t=1$，无需主值。自变量不取复值。
\end{definition}

对实数 $x>1$，
\begin{equation}
  \boxed{
  \Li(x)
  \;=\;
  \li(x)-\li(2)
  }.
  \label{eq:Li-li-real}
\end{equation}
因而
\begin{equation}
  \li(x)\;\sim\;\Li(x)\;\sim\;\frac{x}{\log x}
  \qquad(x\to+\infty).
  \label{eq:li-Li-asymp}
\end{equation}
完整渐近级数见第~\ref{ch:prime_distribution}~章~\eqref{eq:Li-parts-1}--\eqref{eq:Li-asymptotic}。

实比较中 \(\li\) 与 \(\Li\) 可互换，但须按~\eqref{eq:Li-li-real} 调整常数：
含 $\li(x)$ 与 $-\log 2$ 的主项改写为 $\Li$ 时，常数应为 $\li(2)-\log 2$，不可删去 $-\log 2$。
素数定理 $\pi\sim\Li$ 用实函数 $\Li$（第~\ref{ch:prime_number_theorem}~章）。

% ============================================================
\section{\texorpdfstring{$\li$}{li} 的解析延拓}
\label{sec:li-complex-power}
% ============================================================

\subsection{复自变量上的定义}

\begin{definition}{对数积分（复延拓）}{}
\label{def:li-complex}
对 $z\notin[1,\infty)$，取自 $0$ 至 $z$、避开支点 $t=0$ 与 $t=1$ 的路径，
\begin{equation}
  \li(z)
  \;=\;
  \int_{0}^{z}\frac{\mathrm{d}t}{\log t}.
  \label{eq:li-complex}
\end{equation}
$\log$ 取定一分支；路径与分支固定后 $\li(z)$ 单值。
\end{definition}

当 $z=x>1$ 且路径取实轴主值时，~\eqref{eq:li-complex} 与~\eqref{eq:li-def-ch2} 重合。
这是第~\ref{ch:analytic_continuation_general}~章意义下的路径积分延拓：
所得为 $\mathbb{C}\setminus[1,\infty)$ 上固定分支后的单值全纯函数，
而非全平面无割线的整函数；换绕 $t=1$ 的道路一般改变值。

\begin{figure}[htbp]
\centering
\begin{tikzpicture}[
  >=Stealth,
  font=\small,
  scale=1.0
]
  \draw[->] (-0.8,0) -- (6.8,0) node[right] {$\operatorname{Re}t$};
  \draw[->] (0,-1.8) -- (0,3.5) node[above] {$\operatorname{Im}t$};
  \draw[line width=2.2pt, gray!55]
    (1.2,0) -- (6.2,0);
  \node[gray!70!black, below] at (4.2,-0.12) {割线 $[1,+\infty)$};
  \fill[red!70!black] (1.2,0) circle (1.8pt);
  \node[red!70!black, below left] at (1.15,-0.15) {$t=1$};
  \fill (0,0) circle (1.5pt);
  \node[below left] at (0,0) {$0$};
  \draw[blue!75!black, very thick, ->]
    (0.08,0.08)
    .. controls (0.9,1.5) and (2.2,2.4) ..
    (4.5,2.1);
  % 标签与公式分置，避免与路径 $\gamma$ 重叠
  \node[blue!75!black, left] at (1.85,2.05) {路径 $\gamma$};
  \fill[blue!75!black] (4.5,2.1) circle (2pt);
  \node[blue!75!black, right] at (4.65,2.15) {$z\notin[1,+\infty)$};
  \node[anchor=south east, align=right, font=\small, text=black!80]
    at (6.3,2.85) {$\li(z)=\int_{\gamma}\frac{\mathrm{d}t}{\log t}$};
\end{tikzpicture}
\caption{复平面上的 $\li(z)$：路径 $\gamma$ 避开支点 $t=0,1$ 与割线 $[1,+\infty)$
（$\log$ 取定分支）。}
\label{fig:li-complex-path}
\end{figure}

$\li(z)$ 一般取复值。$\Li$ 的自变量限于实数，故复点上使用 $\li(z)$；
没有 $\li(z)=\Li(z)+\li(2)$ 在复域上的对应恒等式。

\subsection{模曲面 \texorpdfstring{$|\li(z)|$}{|li(z)|}}
\label{sec:li-modulus-surface}

图~\ref{fig:li-modulus} 以高度 $|\li(z)|$ 展示定义~\ref{def:li-complex}（金黄色模曲面，三轴过原点）。
奇点、零点与割线在模上分别表现为竖起、触底与沿缝挖空。图为示意。

\begin{figure}[htbp]
\centering
\includegraphics[width=0.98\textwidth]{全书图形/13-li模曲面.pdf}
\caption{$\li$ 的模曲面 $z=|\li(z)|$。
红点：支点 $z=1$；实零点 Soldner 常数 $\mu\approx 1.45137$（$\li(\mu)=0$）。
割线 $[1,+\infty)$ 按主支挖去；近 $z=1$ 处 $|\li|\to+\infty$（高度已裁剪）。}
\label{fig:li-modulus}
\end{figure}

\paragraph{支点 $z=1$ 与实零点 $\mu$}
被积函数 $1/\log t$ 在 $t=1$ 有简单极点；$\li(z)$ 在 $z=1$ 是\textbf{支点}（非 $\Gamma$ 型极点），
$z\to 1$ 时 $|\li(z)|\to+\infty$，割线取 $[1,+\infty)$。
在 $(1,+\infty)$ 上 $\li$ 恰有一个实零点（Soldner 常数）
\begin{equation}
  \li(\mu)=0,
  \qquad
  \mu\approx 1.45136923488\ldots.
  \label{eq:soldner-mu}
\end{equation}
符号：$\li(x)<0$ 当 $1<x<\mu$；$\li(x)>0$ 当 $x>\mu$。
由~\eqref{eq:Li-li-real}，$\Li(2)=0$ 而 $\Li(\mu)=-\li(2)\neq 0$：
$\Li$ 的实零点是 $2$，$\li$ 的是 $\mu$。

\begin{center}
\renewcommand{\arraystretch}{1.25}
\begin{tabular}{lll}
\hline
对象 & 模「竖起」 & 模「触底」 \\
\hline
$\Gamma(s)$ & 单极点 $s=0,-1,\ldots$ & 无 \\
$\zeta(s)$ & $s=1$ 单极点 & 零点 \\
$\li(z)$ & 支点 $z=1$ 及割线 & 实零点 $\mu$ 等 \\
\hline
\end{tabular}
\end{center}

\subsection{大模长时的主项}

在割线平面 $\mathbb{C}\setminus[1,\infty)$ 内，当 $|z|\to\infty$ 且辐角与割线保持正距离时，
\begin{equation}
  \li(z)
  \;\sim\;
  \frac{z}{\log z}.
  \label{eq:li-z-asymp}
\end{equation}
（与实渐近~\eqref{eq:li-asymp-real} 同型；亦可由 $\Ei(\Log z)\sim z/\Log z$ 得，见~\eqref{eq:li-z-Ei}。）
任何满足 $z\notin[1,\infty)$ 的复点均可代入~\eqref{eq:li-complex}；
后文若在特定点上取值，只是~\eqref{eq:li-complex} 的特例，不另定义新函数。

% ============================================================
\section{指数积分 \texorpdfstring{$\Ei$}{Ei}：\texorpdfstring{$\li$}{li} 的变换形式}
\label{sec:Ei}
% ============================================================

\begin{definition}{指数积分}{}
\label{def:Ei}
\[
  \Ei(z)
  =
  -\mathrm{p.v.}\int_{-z}^{+\infty}\frac{e^{-t}}{t}\,\mathrm{d}t,
  \qquad
  z\in\mathbb{C}\setminus(-\infty,0].
\]
\end{definition}

分支一致时，
\begin{equation}
  \li(x)=\Ei(\log x),
  \qquad
  \Li(x)=\Ei(\log x)-\Ei(\log 2)
  \qquad(x>1\text{ 实}),
  \label{eq:li-Ei}
\end{equation}
\begin{equation}
  \li(z)
  \;\longleftrightarrow\;
  \Ei(\Log z)
  \qquad\bigl(z\notin[1,\infty)\bigr).
  \label{eq:li-z-Ei}
\end{equation}
$\Ei$ 是 $\li$ 的变换形式，不引入新的算术对象；理论叙述以 $\li(z)$ 为准。
数值实现见第~\ref{ch:riemann_numerical_approximation}~章。

\subsection{为何引入 \texorpdfstring{$\Ei$}{Ei}}

换元 $t=e^{u}$ 把 $\mathrm{d}t/\log t$ 化为 $e^{u}/u\,\mathrm{d}u$，
$t=1$ 的奇点变为 $u=0$ 处的标准核，归入 $\Ei$ 的主值与分支理论。
$\Ei(w)\sim e^{w}/w$ 直接给出~\eqref{eq:li-z-asymp}。
公式叙述仍写 $\li(z)$；$\Ei$ 不替换其所表示的对象。

% ============================================================
\section{用法分工}
\label{sec:Li-li-convention}
% ============================================================

\begin{center}
\renewcommand{\arraystretch}{1.35}
\begin{tabular}{lll}
\hline
场合 & 所用对象 & 依据 \\
\hline
素数定理、与 $\pi$ 的实比较 &
  实函数 $\Li(x)$ &
  定义~\ref{def:Li-real}；第~\ref{ch:prime_number_theorem}~章 \\
复点上的取值 &
  $\li(z)$ &
  定义~\ref{def:li-complex} \\
数值计算 &
  $\Ei(\Log z)$ 等 &
  \eqref{eq:li-Ei}--\eqref{eq:li-z-Ei} \\
\hline
\end{tabular}
\end{center}

实轴上可用~\eqref{eq:Li-li-real} 在 $\li$ 与 $\Li$ 之间改写；
复点上只能用 $\li(z)$，没有「复自变量的 $\Li$」。

% ============================================================
\section*{本章小结与后续展望}
\addcontentsline{toc}{section}{本章小结与后续展望}
% ============================================================

要点：$\Li=\li-\li(2)$（实）；复函数为 $\li(z)$（割线 $[1,+\infty)$）；
模曲面标明支点 $z=1$ 与零点 $\mu$；$\Ei$ 为变换形式。
下一章在特定复点上应用 $\li(z)$；数值章用 $\Ei$。
